\section{Discussion and Limitations}
\label{sec:limits}

\paragraph{Modelling assumptions.}
The analysis assumes (A1) a target's score on a task can be raised independently of its other scores, (A2) rivals do not react, (A3) training on one task neither helps nor hurts other tasks (no spillover), and (A4) scores lie in $[0,100]$ and gains are limited by the cap, and (A5) investment costs one unit per score point on every task and at every level (linear, task-symmetric cost). Each can fail: training on one task can move others; rivals can respond \citep{chen2026incentives}; ceiling effects make late points dearer: under a quadratic cost the minimum is still cheaper than the mean for \cvxHelmQuadraticShare\% (HELM) and \cvxOllbQuadraticShare\% (Open LLM Leaderboard) of targets (Appendix~\ref{app:sens}), but we did not vary the price across tasks; and our model does not separate legitimate training from contamination \citep{zhou2023cheater,yang2023rephrased}. The caps are assumptions, not predictions.

\paragraph{Scope.}
The targets are all models of two frozen snapshots, many of which are fine-tunes or merges of one another, so rows are dependent and intervals over targets are descriptive. The Open LLM Leaderboard has $n=6$ aggregate tasks and thousands of models; HELM MMLU has 57 subjects of one benchmark, which are strongly correlated, and 79 models. Results do not extrapolate to leaderboards with other sizes or rules. The Condorcet requirement is sufficient for first place under any Condorcet-consistent rule, so its cost is an upper bound for that family; Kemeny and Schulze were not run.

\paragraph{Unfinished, failed and censored analyses.}
(i) The Open LLM Leaderboard sub-tasks (of BBH and MMLU-Pro) were not analysed; HELM's 57 subjects stand in for many-task behaviour. (ii) On the Open LLM Leaderboard the ordinal rules cover every target of ranks 2--100 and, in the main tables and slices, a rank-stratified sample of 101--102 of the 4495 targets; nothing is said about them beyond rank 100, and the minimum cap $h^\ast$ of these rules is computed for ranks 2--100 only. The trimmed mean drops $t=1$ of 6 scores at each end on the Open LLM Leaderboard (all targets) and $t=5$ of 57 on HELM (36 targets), so its two lines are not comparable; 4 of its 144 HELM solves are censored by the time limit. A censored solve is reported as censored, never as infeasible, and one Borda $h^\ast$ row, for a target of rank 4496 in the stratified sample, is censored at the solver limit (it counts neither as feasible nor as infeasible); no other censored row occurs in the main grids or the $h^\ast$ files. (iii) The noise and task-bootstrap experiments use rank-stratified samples of \noiseNHelm{} (HELM) and \noiseNOllb{} (Open LLM Leaderboard, score rules only) targets; they are not population estimates. (iv) Earlier runs with out-of-memory failures, a wrong trimmed-mean formulation or time-outs counted as infeasible were discarded and re-run. (v) The empirical-gain caps rest on 540 fine-tune/base pairs from one leaderboard, not on \emph{targeted} training; on HELM they are transferred from one task.

\paragraph{Ethics and societal impact.}
All data are public aggregate scores; no model is trained and no human subjects or private data are involved. The study can inform benchmark designers about which rules are cheap to move, and can be read as a cost estimate for gaming a leaderboard; no developers are identified.

\paragraph{Use of LLM-based agents.}
Planning, code, figures and this text were produced by LLM-based agents; the solvers use no language model, and every number in the text is generated from the result files by \texttt{paper/make\_numbers.py}.
